\section{Introduction}

Hypersonic thermal protection systems (TPS) play a critical role in safeguarding hypersonic vehicles as they are exposed to extreme aerodynamic heating during flight. Therefore, accurate modeling and prediction of TPS behavior become essential for designing effective thermal protection strategies. Nonetheless, the highly non-linear fluid-thermal-structural interactions, temperature-dependent material property variations, and uncertainties in thermal contact resistances (TCR) at the interfaces of the TPS layers, pose significant challenges for traditional modeling approaches such as finite-element methods (FEM), as well as reduced-order (ROM) and surrogate (SM) modeling techniques.

Particularly, a critical artifact in transient reduced-order modeling of hypersonic TPS is the presence of TCRs. Due to manufacturing limitations, when two solids are brought together to form an interface, the true solid-to-solid contact area between them is generally a small fraction of the apparent area over which they meet {Minges1970}. The net effect of imperfect contact is the formation of a temperature discontinuity at the interface that is in general challenging to predict, and that significantly affects the heat transfer across the interface. The TCRs at the interfaces can vary by orders of magnitude depending on the surface roughness, contact pressure, and temperature {Minges1970,Deng2026}. This variability in TCRs can lead to significant uncertainties in the temperatures predicted by ROMs, and thus compromise the reliability of the ROMs for design and analysis of hypersonic TPS {Fang2023}. Thus, there is a critical need for modeling approaches that can effectively account for the uncertainties in TCRs while still providing accurate predictions of TPS behavior under varying conditions such as boundary conditions and material properties.

Attempts to address the challenges of modeling hypersonic TPS with unknown TCRs have been made using various approaches, ranging from traditional physics-based modeling to data-driven techniques {X}.


The objectives of this work as summarized as follows:
\begin{enumerate}
    \item Demonstrate the systematic application of PIROM to hypersonic TPS with unknown contact resistances, emphasizing on the simultaneous inference of contact conductances and hidden dynamics under varying boundary conditions and material properties.
    \item Benchmark the performance of PIROM against state-of-the-art machine learning, SciML, and traditional physics-based modeling approaches (e.g., FEM) for both inferring contact conductances and transient TPS temperature dynamics.
    \item Assess the performance of the PIROM in an experimental setting,...
\end{enumerate}